\chapter{Lesson 3: Weather}

\section{Vocabulary}

\begin{center}
\begin{tabular}{llll}
\toprule
Characters & Pinyin & Part of speech & Meaning \\
\midrule
\zh{天气} & tiānqì & noun & weather \\
\zh{下雨} & xià yǔ & verb-object & to rain \\
\zh{下雪} & xià xuě & verb-object & to snow \\
\zh{刮风} & guā fēng & verb-object & to be windy \\
\zh{晴天} & qíngtiān & noun & sunny day \\
\zh{阴天} & yīntiān & noun & overcast day \\
\zh{冷} & lěng & adjective & cold \\
\zh{热} & rè & adjective & hot \\
\zh{暖和} & nuǎnhuo & adjective & warm \\
\zh{凉快} & liángkuai & adjective & cool \\
\zh{舒服} & shūfu & adjective & comfortable \\
\zh{闷热} & mēnrè & adjective & hot and humid \\
\zh{预报} & yùbào & noun / verb & forecast \\
\zh{气温} & qìwēn & noun & air temperature \\
\zh{度} & dù & measure word & degree \\
\zh{季节} & jìjié & noun & season \\
\zh{春天} & chūntiān & noun & spring \\
\zh{夏天} & xiàtiān & noun & summer \\
\zh{秋天} & qiūtiān & noun & autumn \\
\zh{冬天} & dōngtiān & noun & winter \\
\zh{比} & bǐ & preposition & compared with \\
\zh{更} & gèng & adverb & even more \\
\zh{还} & hái & adverb & even (more); still \\
\zh{会} & huì & modal & will (likely) \\
\zh{好像} & hǎoxiàng & adverb & it seems \\
\zh{糟糕} & zāogāo & adjective & terrible; ``oh no!'' \\
\zh{要\ldots{}了} & yào \ldots le & & about to \ldots \\
\zh{最} & zuì & adverb & most \\
\bottomrule
\end{tabular}
\end{center}

\section{Comparisons with \zh{比}}

\begin{definition}[The \zh{比} sentence]
A \zh{比} B + adjective (+ difference) ``A is more \emph{adj} than B (by \ldots)''.
\end{definition}

\begin{center}
\begin{tikzpicture}[w/.style={draw, rounded corners, minimum height=0.85cm, align=center, font=\small}]
  \node[w, fill=blue!10] (a) at (0,0) {\zh{今天}\\ A};
  \node[w, fill=red!15] (b) at (1.8,0) {\zh{比}};
  \node[w, fill=blue!10] (c) at (3.6,0) {\zh{昨天}\\ B};
  \node[w, fill=gray!15] (d) at (5.4,0) {(\zh{还/更})\\ adverb};
  \node[w, fill=orange!20] (e) at (7.2,0) {\zh{冷}\\ adjective};
  \node[w, fill=green!15] (f) at (9.3,0) {\zh{五度}\\ difference};
  \foreach \p/\q in {a/b,b/c,c/d,d/e,e/f} \draw[->, thick] (\p) -- (\q);
  \node[font=\small] at (4.65,-1.1) {``Today is (even) five degrees colder than yesterday.''};
\end{tikzpicture}
\end{center}

\begin{itemize}
  \item \zh{北京比滑铁卢热。} Beijing is hotter than Waterloo.
  \item \zh{今天比昨天冷一点儿。} Today is a little colder than yesterday.
  \item \zh{这儿的冬天比我家的冬天冷得多。} / \zh{冷多了。} Winter here is much colder
        than at home.
  \item \zh{明天会比今天更冷。} Tomorrow will be even colder than today.
  \item \zh{我比我妹妹大三岁。} I'm three years older than my younger sister.
\end{itemize}

\begin{important}[What cannot go in a \zh{比} sentence]
\begin{itemize}
  \item No \zh{很}, \zh{非常}, \zh{太} before the adjective: *\zh{今天比昨天很冷}.
        Use \zh{更} or \zh{还} (``even''), or add \zh{得多}/\zh{多了}/\zh{一点儿} after it.
  \item The difference goes \emph{after} the adjective: \zh{高五度}, not *\zh{五度高}.
\end{itemize}
\end{important}

\subsection{Negative comparisons}

\begin{itemize}
  \item A \zh{没有} B (\zh{那么/这么}) adj: ``A is not as \ldots as B''.
        \zh{今天没有昨天那么冷。} Today isn't as cold as yesterday.
  \item A \zh{不比} B adj is a rebuttal: ``A is not more \ldots than B (contrary to what you
        said)''. \zh{上海不比北京冷。}
\end{itemize}

\subsection{Comparisons with verbs}

With a degree complement, \zh{比} B can go before the verb or before the adjective.
\begin{itemize}
  \item \zh{他比我说得快。} / \zh{他说得比我快。} He speaks faster than I do.
\end{itemize}

\subsection{Superlatives with \zh{最}}

\zh{滑铁卢一年中二月最冷。} In Waterloo, February is the coldest month of the year.

\section{Change of state: sentence-final \zh{了}}

In 101R we used \zh{了} after a verb for a completed action. \zh{了} at the \emph{end of
a sentence} announces that a situation is new or has changed.

\begin{center}
\begin{tabular}{lll}
\toprule
Sentence & Meaning & Implied before \\
\midrule
\zh{下雨了。} & It's raining (now). & it wasn't \\
\zh{天气冷了。} & It's gotten cold. & it was warm \\
\zh{我不冷了。} & I'm not cold any more. & I was cold \\
\zh{他二十岁了。} & He's twenty now. & he was younger \\
\zh{春天了。} & It's spring now. & winter \\
\bottomrule
\end{tabular}
\end{center}

\subsection{Imminent action: \zh{要\ldots{}了}, \zh{快(要)\ldots{}了}}

\begin{itemize}
  \item \zh{要下雪了。} It's about to snow.
  \item \zh{快考试了。} Exams are almost here.
  \item \zh{飞机快要起飞了。} The plane is about to take off.
\end{itemize}
With a specific time, use \zh{就要\ldots{}了}: \zh{明天就要放假了。} (not *\zh{快要}).

\section{Probability: \zh{会}}

\zh{会} ``will (probably)'' predicts the future.
\begin{itemize}
  \item \zh{明天会下雨吗？}---\zh{预报说不会。} Will it rain tomorrow?---The forecast says no.
  \item \zh{看样子，晚上会刮大风。} It looks like it'll be very windy tonight.
\end{itemize}

\section{Temperatures}

China uses Celsius (\zh{摄氏} shèshì), like Canada; the US uses Fahrenheit (\zh{华氏}
huáshì). Below zero: \zh{零下十五度} = $-15^\circ$C.

\begin{center}
\begin{tikzpicture}[x=0.75cm, y=0.12cm]
  \draw[->] (0,-20) -- (0,35) node[above, font=\scriptsize] {$^\circ$C};
  \foreach \t in {-20,-10,0,10,20,30} {
    \draw (-0.1,\t) -- (0.1,\t);
    \node[left, font=\scriptsize] at (-0.1,\t) {\t};
  }
  \draw[gray!50] (0,0) -- (13,0);
  \foreach \m/\w/\b in {1/-7/-3, 4/6/14, 7/21/27, 10/9/13} {
    \pgfmathsetmacro{\xa}{\m*1.0}
    \pgfmathsetmacro{\xb}{\m*1.0+0.8}
    \fill[blue!50] (\xa,0) rectangle ++(0.75,\w);
    \fill[red!60] (\xb,0) rectangle ++(0.75,\b);
  }
  \node[font=\scriptsize] at (1.8,-24) {\zh{一月}};
  \node[font=\scriptsize] at (4.8,-24) {\zh{四月}};
  \node[font=\scriptsize] at (7.8,-24) {\zh{七月}};
  \node[font=\scriptsize] at (10.8,-24) {\zh{十月}};
  \fill[blue!50] (11.8,30) rectangle ++(0.4,3); \node[right, font=\scriptsize] at (12.2,31.5) {\zh{滑铁卢}};
  \fill[red!60] (11.8,24) rectangle ++(0.4,3); \node[right, font=\scriptsize] at (12.2,25.5) {\zh{上海}};
\end{tikzpicture}
\end{center}
(Approximate monthly average temperatures, rounded, for illustration.)

\begin{example}[Converting]
$F = \tfrac95 C + 32$. A forecast high of 30$^\circ$C is $\tfrac95(30)+32 = 86^\circ$F;
a Waterloo January night of $-20^\circ$C is $-4^\circ$F.
\end{example}

\begin{aside}
China is so large that it covers almost every climate. Harbin (\zh{哈尔滨}) holds an ice
festival in January at $-25^\circ$C, while Guangzhou (\zh{广州}) is \zh{闷热} for half the
year. In summer southern China has a rainy season called \zh{梅雨} méiyǔ ``plum rain''.
Students from southern China at Waterloo often say Ontario winter is the coldest thing
they have experienced; students from the north find Waterloo's winter comparatively
\zh{暖和}, but the wind brutal.
\end{aside}

\section{Model weather report}

\zh{今天是阴天，气温比昨天低，最高气温零下五度，最低零下十二度。下午会刮风，晚上要下雪了。明天天气会更冷，出门的时候别忘了戴帽子和手套。}

(\zh{出门} chūmén = go out; \zh{忘} wàng = forget.)

\section*{Review}
\begin{itemize}
  \item A \zh{比} B adj (+ difference / \zh{一点儿} / \zh{得多}); \zh{更/还} before the
        adjective; never \zh{很}.
  \item A \zh{没有} B (\zh{那么}) adj for ``not as \ldots as''.
  \item Sentence-final \zh{了} = change of state.
  \item \zh{要/快要/就要\ldots{}了} = about to happen.
  \item \zh{会} for predictions; \zh{最} for superlatives.
  \item Seasons \zh{春夏秋冬}; \zh{零下} for below zero.
\end{itemize}
